% SEGMENT OLP-0346-S01
% Part: lambda-calculus
% Chapter: introduction
% Section: church-rosser

\documentclass[../../../include/open-logic-section]{subfiles}

\begin{document}

\olfileid{lam}{int}{cr}
\olsection{சர்ச்--ரோசர் பண்பு}

\begin{thm}
\ollabel{thm:church-rosser}
$M$, $N_1$, $N_2$ ஆகியவை $M \red N_1$, $M \red N_2$ என அமையும் சொற்களாக
இருக்கட்டும். அப்போது $N_1 \red P$, $N_2 \red P$ என அமையும் சொல் $P$ ஒன்று
உள்ளது.
\end{thm}

\begin{cor}
$M$ ஐ இயல்புவடிவத்திற்குக் குறைக்க முடியும் என்று கொள்க. அப்போது இந்த
இயல்புவடிவம் தனித்தது.
\end{cor}

\begin{proof}
$M \red N_1$, $M \red N_2$ எனில், முந்தைய தேற்றப்படி $N_1$, $N_2$ இரண்டும்
~$P$ ஆகக் குறையும் சொல்~$P$ ஒன்று உள்ளது. $N_1$,~$N_2$ இரண்டும்
இயல்புவடிவத்தில் இருந்தால், $N_1 \ident P \ident N_2$ எனில் மட்டுமே இது
நிகழ முடியும்.
\end{proof}

இறுதியாக, இரு சொற்கள் $M$,~$N$ ஒரு பொதுவான சொல்லாகக் குறைந்தால் அவை
\emph{$\beta$-சமானமானவை}, அல்லது வெறுமனே \emph{சமானமானவை}, என்று கூறுவோம்;
வேறு சொற்களில், $M \red P$, $N \red P$ என அமையும் ஏதோ $P$ இருந்தால் இவ்வாறு
கூறுவோம். இதை $M \equal[\beta] N$ என்று எழுதுகிறோம்.
\olref{thm:church-rosser} ஐப் பயன்படுத்தி, $\equal[\beta]$ ஒரு சமானத் தொடர்பு
என்பதையும், ஒவ்வொரு $M$,~$N$ க்கும் $M \red N$ அல்லது $N \red M$ எனில்
$M \equal[\beta] N$ என்ற கூடுதல் பண்பு அதற்கு உள்ளதையும் நீங்கள் சரிபார்க்கலாம்.
(உண்மையில், இப்பண்பைக் கொண்ட \emph{மீச்சிறு} சமானத் தொடர்பு
$\equal[\beta]$ என்று காட்டலாம்.)

\end{document}
